\chapter{OPTIMIZATION (CORE OF HYBRID ALGORITHMS)}

This block focuses on classical optimization methods that are essential for hybrid quantum-classical algorithms.

\begin{center}
\textbf{Variational quantum algorithms fundamentally rely on classical optimization to adjust parameters and minimize objective functions}
\end{center}

\vspace{0.5cm}

\section*{1. FUNDAMENTAL PROBLEM}

In hybrid algorithms such as VQE and QAOA, the goal is to minimize an objective function:

\[
\min_{\theta} f(\theta) = \langle \psi(\theta) | H | \psi(\theta) \rangle
\]

where:
\begin{itemize}
\item $\theta$ represents circuit parameters
\item $H$ is the Hamiltonian or cost function
\item $f(\theta)$ is evaluated using a quantum circuit
\end{itemize}

This creates a feedback loop:
\begin{enumerate}
\item Prepare quantum state $|\psi(\theta)\rangle$
\item Measure expectation value
\item Update parameters using a classical optimizer
\end{enumerate}

\vspace{0.5cm}

\section*{2. CLASSICAL OPTIMIZATION METHODS}

\subsection*{2.1 Gradient Descent}

A first-order optimization method based on the gradient of the objective function.

Update rule:

\[
\theta_{k+1} = \theta_k - \eta \nabla f(\theta_k)
\]

where:
\begin{itemize}
\item $\eta$ is the learning rate
\item $\nabla f(\theta)$ is the gradient
\end{itemize}

\subsubsection*{Key Properties}

\begin{itemize}
\item Simple and widely used
\item Requires gradient computation
\item Sensitive to learning rate choice
\end{itemize}

\vspace{0.3cm}

\subsection*{2.2 Adam (Adaptive Moment Estimation)}

An adaptive optimization algorithm combining momentum and RMSProp ideas.

Update equations:

\[
m_t = \beta_1 m_{t-1} + (1 - \beta_1)\nabla f(\theta_t)
\]

\[
v_t = \beta_2 v_{t-1} + (1 - \beta_2)(\nabla f(\theta_t))^2
\]

\[
\theta_{t+1} = \theta_t - \eta \frac{m_t}{\sqrt{v_t} + \epsilon}
\]

\subsubsection*{Key Properties}

\begin{itemize}
\item Adaptive learning rates
\item Faster convergence in practice
\item Robust for noisy gradients
\end{itemize}

\vspace{0.3cm}

\subsection*{2.3 COBYLA (Constrained Optimization BY Linear Approximation)}

A gradient-free optimization method.

\subsubsection*{Core Idea}

\begin{itemize}
\item Approximates objective function locally using linear models
\item Handles constraints naturally
\item Does not require gradient information
\end{itemize}

\subsubsection*{Applications}

\begin{itemize}
\item Variational quantum algorithms with noisy measurements
\item Optimization under constraints
\end{itemize}

\vspace{0.3cm}

\subsection*{2.4 SPSA (Simultaneous Perturbation Stochastic Approximation)}

A stochastic gradient approximation method particularly suited for quantum systems.

Gradient approximation:

\[
\hat{\nabla} f(\theta) \approx \frac{f(\theta + \Delta) - f(\theta - \Delta)}{2\Delta}
\]

where:
\begin{itemize}
\item $\Delta$ is a random perturbation vector
\end{itemize}

\subsubsection*{Key Properties}

\begin{itemize}
\item Requires only two function evaluations per iteration
\item Efficient for high-dimensional parameter spaces
\item Robust to noise
\end{itemize}

\vspace{0.5cm}

\section*{3. COMPARISON OF METHODS}

\begin{center}
\begin{tabular}{|c|c|c|c|}
\hline
Method & Gradient Required & Noise Robustness & Efficiency \\
\hline
Gradient Descent & Yes & Low & Medium \\
Adam & Yes & High & High \\
COBYLA & No & Medium & Medium \\
SPSA & Approximate & Very High & Very High \\
\hline
\end{tabular}
\end{center}

\vspace{0.5cm}

\section*{4. ROLE IN QUANTUM ALGORITHMS}

In hybrid quantum algorithms:

\begin{itemize}
\item The quantum computer evaluates $f(\theta)$
\item The classical optimizer updates $\theta$
\item The process iterates until convergence
\end{itemize}

This hybrid loop can be written as:

\[
\theta \rightarrow \text{Quantum Circuit} \rightarrow f(\theta) \rightarrow \text{Optimizer} \rightarrow \theta'
\]

\vspace{0.5cm}

\section*{5. IMPLEMENTATION EXAMPLE}

\begin{minted}{python}
import numpy as np

def f(theta):
    return theta[0]**2 + theta[1]**2

theta = np.array([1.0, 1.0])
eta = 0.1

for _ in range(100):
    grad = np.array([2*theta[0], 2*theta[1]])
    theta = theta - eta * grad

print(theta)
\end{minted}

\vspace{0.5cm}

\section*{FINAL CONNECTION}

Optimization is the central bridge between:

\begin{enumerate}
\item Quantum circuits (state preparation)
\item Measurement (objective evaluation)
\item Classical computation (parameter update)
\end{enumerate}

Without efficient optimization, variational quantum algorithms cannot function effectively.

\vspace{0.5cm}

\section*{STUDY ROADMAP}

\subsection*{Phase 1 --- Fundamentals}
\begin{itemize}
\item Objective functions
\item Gradients and derivatives
\end{itemize}

\subsection*{Phase 2 --- Gradient-Based Methods}
\begin{itemize}
\item Gradient descent
\item Adam optimizer
\end{itemize}

\subsection*{Phase 3 --- Gradient-Free Methods}
\begin{itemize}
\item COBYLA
\item SPSA
\end{itemize}

\subsection*{Phase 4 --- Applications}
\begin{itemize}
\item VQE optimization
\item QAOA parameter tuning
\end{itemize}